\documentclass[aspectratio=169,11pt]{beamer}
\makeatletter
\input{../../templates/dune-professional/beamerthemeDUNEProfessional.sty}
\makeatother
\title{DPS hardware for vertical drift}
\subtitle{Existing BDE design, additions and procurement}
\author{Manuel Arroyave}
\date{I\&I workshop draft --- 27 October 2026}
\newcommand{\mydate}{27 October 2026}
\hypersetup{colorlinks=true,urlcolor=DuneInfo,linkcolor=DuneInfo}
\newcommand{\planningnote}[1]{\par\medskip{\footnotesize\color{DuneTextMuted}#1\par}}
\begin{document}

\begin{frame}[plain]
  \titlepage
\end{frame}

\begin{frame}{Use a commercial PLC for VD detector protection}
  \large
  \textbf{Proposal:} use the bottom drift electronics (BDE) design
  and add the inputs and outputs needed by the other VD systems.\par
  \bigskip
  Put the central programmable logic controller (PLC) in the CUC,
  with distributed I/O near the equipment.\par
  \bigskip
  OPC UA carries DPS status and supervisory requests to and from slow controls.\par
  \planningnote{CUC = Central Utility Cavern; central-PLC location proposed. HD can later use an extension or a separate PLC. This plan covers VD.}
\end{frame}

\begin{frame}{The VD rack plan already reserves space for DPS}
  \renewcommand{\arraystretch}{1.4}
  \begin{tabularx}{\textwidth}{@{}>{\RaggedRight\arraybackslash}p{3.4cm}r>{\RaggedRight\arraybackslash}X@{}}
    \toprule
    \textbf{Rack type} & \textbf{Count} & \textbf{DPS equipment / connections}\\
    \midrule
    BDE auxiliary & 10 & 12U per rack; eight WIEC connections per rack\\
    BDE low voltage & 20 & One PL506 per rack; four WIECs per supply\\
    Bias high voltage & 10 & 12U per rack for the DPS interface to ISEG supplies\\
    \bottomrule
  \end{tabularx}\par
  \bigskip
  The auxiliary racks contain the WIEC I/O; the bias racks contain the ISEG inhibit outputs.
  \planningnote{May 2025 VD rack workbook; June 2026 BDE layout. U is rack height. These are design allocations; the bias-HV sheet is a first estimate.}
\end{frame}

\begin{frame}{One BDE I/O group serves eight WIECs}
  \centering
  \begin{tikzpicture}[
    box/.style={fill=DunePaper,text=DuneOnAccent,align=center,
      text width=3.1cm,minimum height=1.05cm,inner sep=6pt},
    edge/.style={-{Stealth},draw=DuneSky,line width=1.2pt}]
    \node[box,fill=DuneSky] (io) at (-4,0) {BDE auxiliary rack\\remote I/O};
    \node[box] (lv) at (1.6,1.1) {2 LV racks\\2 PL506 supplies};
    \node[box,fill=DuneCoral] (wiec) at (1.6,-1.1) {8 WIECs\\8 PTCs + 48 WIBs};
    \draw[edge] (io.north east) -| (-1.7,1.1) -- node[above,font=\small] {inhibits} (lv.west);
    \draw[edge,<->] (io.south east) -| (-1.7,-1.1) -- node[below,font=\small] {EtherCAT} (wiec.west);
    \draw[edge] (lv.east) -- (4.35,1.1) -- node[right,font=\small] {48 V} (4.35,-1.1) -- (wiec.east);
  \end{tikzpicture}\par
  \medskip
  \DuneDiagramCaption{Ten groups: 20 PL506 supplies, 80 WIECs and 480 WIBs.}
  \medskip
  \raggedright The auxiliary rack also connects the fans, heaters and temperature sensors.
  \planningnote{WIEC = warm interface electronics crate; PTC = power and timing card; WIB = warm interface board. June 2026 layout; fan/heater wiring omitted here.}
\end{frame}

\begin{frame}{Panel drawings and the Ignition front end are available}
  \begin{columns}[T,onlytextwidth]
    \begin{column}{0.47\textwidth}
      \textbf{TPC protection panel designs}\par\medskip
      \colorbox{DunePaper}{\includegraphics[page=2,width=\dimexpr\linewidth-2\fboxsep\relax,height=3.1cm,keepaspectratio,
        trim=290bp 500bp 900bp 430bp,clip]{figures/tpc-panel-reference.pdf}}\par
      \medskip\DuneDiagramCaption{Central PLC, bias I/O and LV/fan/heater panels.}
    \end{column}
    \begin{column}{0.49\textwidth}
      \textbf{BDE Ignition screens}\par\medskip
      \includegraphics[page=7,width=\linewidth,height=3.1cm,keepaspectratio,
        trim=10bp 210bp 460bp 74bp,clip]{figures/sc-progress-june2026.pdf}\par
      \planningnote{Existing display of WIEC groups; June 2026 SC progress presentation.}
    \end{column}
  \end{columns}\par
  \medskip
  Parts list: CX5140 PLC, EK1100 couplers, EK1521 fibre links and EL2124 inhibit outputs.
  \planningnote{TPC reference: EDMS 2401090 v4, sheet 2; component names from the 2023 parts list.}
\end{frame}

\begin{frame}{The VD bias-panel draft specifies two DIN rails per rack}
  \renewcommand{\arraystretch}{1.3}
  \begin{tabularx}{\textwidth}{@{}>{\RaggedRight\arraybackslash}Xrr@{}}
    \toprule
    \textbf{Component} & \textbf{Per rail} & \textbf{10 racks}\\
    \midrule
    EK1100 EtherCAT coupler & 1 & 20\\
    EL9410 supply terminal & 1 & 20\\
    EL2124: four 5 V outputs & 18 & 360\\
    EL9011 end cap & 1 & 20\\
    24 V power supply & 1 & 20\\
    \bottomrule
  \end{tabularx}\par
  \medskip
  Each rail occupies 337 mm. The draft provides 1,440 output channels;
  the working VD inventory lists 576 bias channels.\par
  \planningnote{May 2025 workbook, Bias HV Rack!K17:K21. Quantities describe the first-pass panel layout.}
\end{frame}

\begin{frame}{VD power estimates for the BDE racks}
  \renewcommand{\arraystretch}{1.4}
  \begin{tabularx}{\textwidth}{@{}>{\RaggedRight\arraybackslash}p{3.7cm}>{\RaggedRight\arraybackslash}p{2.4cm}>{\RaggedRight\arraybackslash}X@{}}
    \toprule
    \textbf{Equipment} & \textbf{Per rack} & \textbf{VD basis}\\
    \midrule
    PL506 input & 2.25 kW & 20 racks: about 45.1 kW for the WIEC supplies\\
    Auxiliary DPS input & About 1.11 kW & 10 racks: about 11.1 kW; eight WIECs per rack\\
    Bias DPS supply allowance & 24 V, 40 A & 960 W DC capacity per rack in the workbook\\
    \bottomrule
  \end{tabularx}\par
  \bigskip
  The auxiliary estimate scales the recorded 831 W from six to eight WIECs:
  $831\times8/6 = 1{,}108$ W.\par
  \planningnote{Planning estimates from the May 2025 workbook. Linear scaling follows its VD note; it is not a measured load. Supply capacity is separate from AC demand.}
\end{frame}

\begin{frame}{The PX3 rack-protection setup already defines shutdown}
  \renewcommand{\arraystretch}{1.45}
  \begin{tabularx}{\textwidth}{@{}>{\RaggedRight\arraybackslash}p{3.7cm}>{\RaggedRight\arraybackslash}X@{}}
    \toprule
    \textbf{June 2026 setup guide} & \textbf{Documented behaviour}\\
    \midrule
    Ignition connected & Ignition reads sensors and switches outlets over SNMP\\
    Ignition unreachable & PX3 local rules handle smoke, water and temperature; the network-switch outlet stays on\\
    After a fallback trip & Power is restored manually or through Ignition\\
    \bottomrule
  \end{tabularx}\par
  \bigskip
  This setup combines Ignition control with autonomous PDU fallback. The VD integration includes the PLC protection permits.
  \planningnote{DUNE-doc-32655 v8, pp.~3--6 and 14--16. The PDU fallback is documented; the PLC permit integration is the proposed extension.}
\end{frame}

\begin{frame}{The VD inventory covers the other subsystem connections}
  \renewcommand{\arraystretch}{1.35}
  \begin{tabularx}{\textwidth}{@{}>{\RaggedRight\arraybackslash}p{3.2cm}>{\RaggedRight\arraybackslash}X@{}}
    \toprule
    \textbf{Subsystem} & \textbf{Working inventory and interfaces}\\
    \midrule
    PDS readout / calibration & 42 DAPHNE boards and 8 calibration boxes; 58 LV feeds from MPOD modules\\
    Power over fibre (PoF) & 80 laser boxes; 640 temperature and 640 voltage readings; compact PLC interface proposed\\
    TDE & 20 mezzanine racks; power supplies and readout-crate management\\
    Calibration / cryogenic instrumentation & Level and temperature readout, purity monitors, ionization laser and neutron source\\
    \bottomrule
  \end{tabularx}
  \planningnote{VD hardware lists and SC architecture. Device monitoring reaches Ignition through native drivers and OPC UA; protection uses the subsystem inhibit interfaces.}
\end{frame}

\begin{frame}{Protection timing includes detection and equipment shutdown}
  \large
  No critical protection function requiring a response below 1 ms
  has been identified in the current review.\par
  \bigskip
  Total response time includes the sensor, communication,
  PLC cycle, output and equipment shutdown.\par
  \bigskip
  The simulation exercises interface load and fault cases.
  Bench tests measure the complete hardware path.\par
  \planningnote{The existing software benchmarks do not measure a PLC trip. Subsystem timing limits and the measured response form the acceptance criteria for each test.}
\end{frame}

\begin{frame}{DPS protection and SC monitoring have defined paths}
  \renewcommand{\arraystretch}{1.35}
  \begin{tabularx}{\textwidth}{@{}>{\RaggedRight\arraybackslash}p{3.5cm}>{\RaggedRight\arraybackslash}X@{}}
    \toprule
    \textbf{Connection} & \textbf{Path in the design}\\
    \midrule
    PLC to detector equipment & EtherCAT remote I/O; PTC fibre links; wired supply inhibits\\
    DPS to slow controls & OPC UA status, trip causes and supervisory requests; Ignition display and history\\
    PDUs and power supplies & SNMP to the SC drivers; local device protection and DPS inhibit interfaces\\
    Facility systems & Cryogenic controls via OPC UA; Metasys via BACnet in the SC drawing; limited SDSTA hardware exchanges proposed\\
    \bottomrule
  \end{tabularx}
  \planningnote{BDE design, supplied SC architecture and the I\&I correspondence. Proposed control rule: an SC operating request cannot override a protection trip.}
\end{frame}

\begin{frame}{The BMS drawing identifies 20 control-panel symbols}
  \renewcommand{\arraystretch}{1.4}
  \begin{tabularx}{\textwidth}{@{}>{\RaggedRight\arraybackslash}p{3.4cm}>{\RaggedRight\arraybackslash}X@{}}
    \toprule
    \textbf{Drawing area} & \textbf{Interfaces on the inspection maps}\\
    \midrule
    CUC, level 4850 & Eight panel references: air handling, chilled water, power and associated equipment\\
    Level 4910 & North and south sump interfaces, SP-1 and SP-2\\
    4850-08-ES group & Spare panel for future equipment; mapped to the access-drift inspection area\\
    Other 4850 areas & Chillers, pumps, ventilation and generator interfaces\\
    \bottomrule
  \end{tabularx}\par
  \medskip
  Inspection planned with Rachel; Dale Curran is the SDSTA integration contact identified for this work.
  \planningnote{U-FD-M-901 rev12, 14 February 2025; marked 4850/4910 maps and survey register. Symbols are drawing references; field observations are pending.}
\end{frame}

\begin{frame}{The VD bench tests use the BDE and G271 equipment}
  \renewcommand{\arraystretch}{1.4}
  \begin{tabularx}{\textwidth}{@{}>{\RaggedRight\arraybackslash}X>{\RaggedRight\arraybackslash}X@{}}
    \toprule
    \textbf{Test} & \textbf{Measure or demonstrate}\\
    \midrule
    PLC, remote I/O, PTC and supply inhibit & Complete trip time and equipment feedback\\
    PDU sensors and shutdown logic & Response with SC connected, lost or stalled\\
    SC commands, reset and re-enable & Protection takes priority; recovery is deliberate\\
    Fieldbus, network and power failures & Output states, lost branches and restart behaviour\\
    \bottomrule
  \end{tabularx}
  \planningnote{July G271 inventory: PDU, PL506, mini-MPOD and network switch installed; WIEC/PLC additions listed separately. The BDE vertical-slice work supplies the PTC/WIB context.}
\end{frame}

\begin{frame}{VD hardware work through the June 2027 review}
  \renewcommand{\arraystretch}{1.4}
  \begin{tabularx}{\textwidth}{@{}>{\RaggedRight\arraybackslash}p{3.3cm}>{\RaggedRight\arraybackslash}X@{}}
    \toprule
    \textbf{Work} & \textbf{Deliverable}\\
    \midrule
    VD hardware adaptation & Panel and module allocation for the additional subsystem signals; CUC controller and fieldbus connections\\
    Prototype tests & Trip-time measurements, fault response, SC/PDU integration and recovery results\\
    June 2027 review & Drawings, parts and cable lists, power budget, costs, lead times and spares\\
    Procurement / installation & Orders after readiness review; delivery and installation sequence for the site\\
    \bottomrule
  \end{tabularx}
  \planningnote{Schedule basis: possible facility interfaces in summer 2027; full VD response-matrix review in early 2028. Existing documents support the hardware work now.}
\end{frame}

\begin{frame}{Hardware work and responsibilities}
  \renewcommand{\arraystretch}{1.35}
  \begin{tabularx}{\textwidth}{@{}>{\RaggedRight\arraybackslash}p{4.2cm}>{\RaggedRight\arraybackslash}X@{}}
    \toprule
    \textbf{Proposed team} & \textbf{Work}\\
    \midrule
    Manuel + control engineer & VD hardware design, SURF interfaces and procurement package\\
    Terri, Trevor + Manuel & BDE panels, PLC/I/O implementation and bench tests\\
    Manuel + consortia & Equipment limits, subsystem signals and protective actions\\
    Manuel + Geoff & SC/DPS command priority, PDU integration and recovery\\
    Vivian + subsystem / DPS teams & ICD updates and interface records\\
    \bottomrule
  \end{tabularx}
  \planningnote{Proposed division from the I\&I discussion. Before DPS is available, Manuel proposes that each consortium remains responsible for protecting its equipment.}
\end{frame}

\appendix
\begin{frame}{Reference: the May 2025 VD rack inventory}
  \renewcommand{\arraystretch}{1.15}
  \begin{tabularx}{\textwidth}{@{}>{\RaggedRight\arraybackslash}Xr>{\RaggedRight\arraybackslash}Xr@{}}
    \toprule
    \textbf{Rack type} & \textbf{Count} & \textbf{Rack type} & \textbf{Count}\\
    \midrule
    BDE low voltage & 20 & PDS low voltage & 4\\
    BDE auxiliary & 10 & Drift HV & 2\\
    TDE & 20 & Reserved drift HV & 2\\
    Bias HV & 10 & Spare / CalCI & 8\\
    \bottomrule
  \end{tabularx}\par
  \bigskip
  \textbf{76 numbered mezzanine positions:} SEDM-01 to SEDM-76.\par
  \medskip
  The worksheet also has a separate concept for 40 top-plate mini-racks.
  \planningnote{May 2025 allocation: rack positions, distinct from PLC panels and remote-I/O stations. Includes the listed reserved and spare positions.}
\end{frame}
\end{document}
